
\documentclass[a4paper,11pt]{article}
\usepackage[a4paper, margin=8em]{geometry}

% usa i pacchetti per la scrittura in italiano
\usepackage[french,italian]{babel}
\usepackage[T1]{fontenc}
\usepackage[utf8]{inputenc}
\frenchspacing 

% usa i pacchetti per la formattazione matematica
\usepackage{amsmath, amssymb, amsthm, amsfonts}

% usa altri pacchetti
\usepackage{gensymb}
\usepackage{hyperref}
\usepackage{standalone}

% imposta il titolo
\title{Appunti Elettronica Digitale}
\author{Luca Seggiani}
\date{2026}

% imposta lo stile
% usa helvetica
\usepackage[scaled]{helvet}
% usa palatino
\usepackage{palatino}
% usa un font monospazio guardabile
\usepackage{lmodern}

\renewcommand{\rmdefault}{ppl}
\renewcommand{\sfdefault}{phv}
\renewcommand{\ttdefault}{lmtt}

% disponi il titolo
\makeatletter
\renewcommand{\maketitle} {
	\begin{center} 
		\begin{minipage}[t]{.8\textwidth}
			\textsf{\huge\bfseries \@title} 
		\end{minipage}%
		\begin{minipage}[t]{.2\textwidth}
			\raggedleft \vspace{-1.65em}
			\textsf{\small \@author} \vfill
			\textsf{\small \@date}
		\end{minipage}
		\par
	\end{center}

	\thispagestyle{empty}
	\pagestyle{fancy}
}
\makeatother

% disponi teoremi
\usepackage{tcolorbox}
\newtcolorbox[auto counter, number within=section]{theorem}[2][]{%
	colback=blue!10, 
	colframe=blue!40!black, 
	sharp corners=northwest,
	fonttitle=\sffamily\bfseries, 
	title=~\thetcbcounter: #2, 
	#1
}

% disponi definizioni
\newtcolorbox[auto counter, number within=section]{definition}[2][]{%
	colback=red!10,
	colframe=red!40!black,
	sharp corners=northwest,
	fonttitle=\sffamily\bfseries,
	title=~\thetcbcounter: #2,
	#1
}

% U.D.M
\newcommand{\amp}{\ensuremath{\mathrm{A}}}
\newcommand{\volt}{\ensuremath{\mathrm{V}}}
\newcommand{\meter}{\ensuremath{\mathrm{m}}}
\newcommand{\second}{\ensuremath{\mathrm{s}}}
\newcommand{\farad}{\ensuremath{\mathrm{F}}}
\newcommand{\henry}{\ensuremath{\mathrm{H}}}
\newcommand{\siemens}{\ensuremath{\mathrm{S}}}

% circuiti
\usepackage{circuitikz}
\usetikzlibrary{babel}

% disegni
\usepackage{pgfplots}
\pgfplotsset{width=10cm,compat=1.9}

% disponi codice
\usepackage{listings}
\usepackage[table]{xcolor}

\definecolor{codegreen}{rgb}{0,0.6,0}
\definecolor{codegray}{rgb}{0.5,0.5,0.5}
\definecolor{codepurple}{rgb}{0.58,0,0.82}
\definecolor{backcolour}{rgb}{0.95,0.95,0.92}

\lstdefinestyle{codestyle}{
		backgroundcolor=\color{black!5}, 
		commentstyle=\color{codegreen},
		keywordstyle=\bfseries\color{magenta},
		numberstyle=\sffamily\tiny\color{black!60},
		stringstyle=\color{green!50!black},
		basicstyle=\ttfamily\footnotesize,
		breakatwhitespace=false,         
		breaklines=true,                 
		captionpos=b,                    
		keepspaces=true,                 
		numbers=left,                    
		numbersep=5pt,                  
		showspaces=false,                
		showstringspaces=false,
		showtabs=false,                  
		tabsize=2
}

\lstdefinestyle{shellstyle}{
		backgroundcolor=\color{black!5}, 
		basicstyle=\ttfamily\footnotesize\color{black}, 
		commentstyle=\color{black}, 
		keywordstyle=\color{black},
		numberstyle=\color{black!5},
		stringstyle=\color{black}, 
		showspaces=false,
		showstringspaces=false, 
		showtabs=false, 
		tabsize=2, 
		numbers=none, 
		breaklines=true
}

\lstdefinelanguage{javascript}{
	keywords={typeof, new, true, false, catch, function, return, null, catch, switch, var, if, in, while, do, else, case, break},
	keywordstyle=\color{blue}\bfseries,
	ndkeywords={class, export, boolean, throw, implements, import, this},
	ndkeywordstyle=\color{darkgray}\bfseries,
	identifierstyle=\color{black},
	sensitive=false,
	comment=[l]{//},
	morecomment=[s]{/*}{*/},
	commentstyle=\color{purple}\ttfamily,
	stringstyle=\color{red}\ttfamily,
	morestring=[b]',
	morestring=[b]"
}

% disponi sezioni
\usepackage{titlesec}

\titleformat{\section}
	{\sffamily\Large\bfseries} 
	{\thesection}{1em}{} 
\titleformat{\subsection}
	{\sffamily\large\bfseries}   
	{\thesubsection}{1em}{} 
\titleformat{\subsubsection}
	{\sffamily\normalsize\bfseries} 
	{\thesubsubsection}{1em}{}

% disponi alberi
\usepackage{forest}

\forestset{
	rectstyle/.style={
		for tree={rectangle,draw,font=\large\sffamily}
	},
	roundstyle/.style={
		for tree={circle,draw,font=\large}
	}
}

% disponi algoritmi
\usepackage{algorithm}
\usepackage{algorithmic}
\makeatletter
\renewcommand{\ALG@name}{Algoritmo}
\makeatother

% disponi numeri di pagina
\usepackage{fancyhdr}
\fancyhf{} 
\fancyfoot[L]{\sffamily{\thepage}}

\makeatletter
\fancyhead[L]{\raisebox{1ex}[0pt][0pt]{\sffamily{\@title \ \@date}}} 
\fancyhead[R]{\raisebox{1ex}[0pt][0pt]{\sffamily{\@author}}}
\makeatother

\begin{document}
% sezione (data)
\section{Lezione del 08-05-26}

% stili pagina
\thispagestyle{empty}
\pagestyle{fancy}

% testo
\subsection{Considerazioni di potenza dell'inverter CMOS}
Al termine della scorsa lezione abbiamo studiato la curva \textbf{VTC} (\textit{Voltage Transfer Characteristic}) della porta NOT in tecnologia CMOS. Facciamo adesso delle considerazioni riguardo alla potenza, statica e dinamica, dissipata.

\begin{enumerate}
\item \textbf{\textsf{Potenza statica}} \\
	Dall'analiis del circuito notiamo che, in ambedue gli stati stabili (uscita $V_{out}$ alta o bassa) non c'è conduzione di corrente in alcun punto del circuito. Questo significherà che, idealmente, la potenza statica dissipata è nulla.

	Notiamo che comunque si avranno \textit{correnti parassite}, cioè i transistori MOSFET avranno una piccola resistenza associata, e la potenza dissipata non sarà nulla;

\item \textbf{\textsf{Potenza dinamica}} \\
	Questa è la potenza dissipata durante i transitori da alto a basso, e da basso ad alto, dell'uscita. Rifacendoci alla caratteristica ricavata alla scorsa lezione, abbiamo che questa è dissipata nella zona in cui entrambi i transistor conducono:
$$
V_{TN} \leq V_{in} \leq V_{DD} + V_{TP}
$$
Chiamiamo questa zona impropriamente "cortocircuito" (non c'è un vero cortocircuito ma solo i transistor in conduzione), e quindi la potenza dissipata \textit{potenza di cortocircuito}.

C'è un altra potenza associata ai transitori, data dalla \textit{carica/scarica capacità}. Questa viene dal fatto che sul nodo di uscita di una porta CMOS sarà sempre presente una certa \textit{capacità equivalente}. La potenza dissipata sarà quella necessaria a caricare la capacità.
\end{enumerate}

\subsubsection{Capacità equivalente dell'inverter CMOS}
Vediamo nel dettaglio queste capacità. Iniziamo dal tracciare uno schema con due porte not in CMOS, in serie:

\begin{center}
\begin{tikzpicture}[transform shape]
	% Paths, nodes and wires:
	\node[pmos, bulk] at (5, 5.77){};
	\node[nmos, arrowmos, bulk] at (5, 3.23){};
	\draw (4.02, 5.77) -| (3.5, 5.75);
	\draw (4.02, 3.23) -| (3.5, 3.25);
	\draw (3.5, 3.27) -| (3.5, 5.77);
	\draw (5, 5) -| (5, 4);
	\draw (5, 5.77) -| (5.25, 6);
	\draw (5.25, 6) |- (5, 6.25);
	\draw (5, 3.23) -| (5.25, 3);
	\draw (5.25, 3) |- (5, 2.75);
	\node[vcc] at (5, 6.54){};
	\node[ground] at (5, 2.46){};
	\node[pmos, bulk] at (9.5, 5.77){};
	\node[nmos, arrowmos, bulk] at (9.5, 3.23){};
	\draw (8.52, 5.77) -| (8, 5.75);
	\draw (8.52, 3.23) -| (8, 3.25);
	\draw (8, 3.25) -| (8, 5.75);
	\draw (9.5, 5) -| (9.5, 4);
	\draw (9.5, 5.77) -| (9.75, 6);
	\draw (9.75, 6) |- (9.5, 6.25);
	\draw (9.5, 3.23) -| (9.75, 3);
	\draw (9.75, 3) |- (9.5, 2.75);
	\node[vcc] at (9.5, 6.54){};
	\node[ground] at (9.5, 2.46){};
	\draw (8, 4.5) -- (5, 4.5);
	\draw (6.5, 4.5) to[capacitor, l={$C_W$}] (6.5, 2.5);
	\node[ground] at (6.5, 2.5){};
	\draw[-stealth] (9.5, 4.5) -- (10.5, 4.5);
	\draw (3.5, 4.5) -- (2.5, 4.5);
	\node[shape=rectangle, minimum width=0.465cm, minimum height=0.465cm] at (2.25, 4.5){} node[anchor=east, align=right, text width=0.077cm, inner sep=6pt] at (2.5, 4.5){$V_{in}$};
	\node[shape=rectangle, minimum width=0.465cm, minimum height=0.465cm] at (10.75, 4.5){} node[anchor=west, align=left, text width=0.077cm, inner sep=6pt] at (10.5, 4.5){$V_{out}$};
	\draw (5.75, 3) to[capacitor, /tikz/circuitikz/bipoles/length=0.700cm, l={$C_{DB}$}] (5.25, 3);
	\draw (5.75, 3) -| (5.75, 4.5);
	\draw (5.25, 6) to[capacitor, /tikz/circuitikz/bipoles/length=0.700cm, l={$C_{DB}$}] (5.75, 6);
	\draw (5.75, 6) -- (5.75, 4.5);
	\draw (3.75, 4.75) to[capacitor, /tikz/circuitikz/bipoles/length=0.700cm, l={$C_{GD}$}] (4.25, 4.75);
	\draw (3.75, 4.25) to[capacitor, /tikz/circuitikz/bipoles/length=0.700cm, l_={$C_{GD}$}] (4.25, 4.25);
	\draw (7.75, 4.5) -- (7.75, 2.5) -- (8, 2.5);
	\draw (8.5, 2.5) -- (9.5, 2.5);
	\draw (7.75, 4.5) -- (7.75, 6.5) -- (8, 6.5);
	\draw (8.5, 6.5) -- (9.5, 6.5);
	\draw (8, 6.5) to[capacitor, /tikz/circuitikz/bipoles/length=0.700cm, l={$C_{GP}$}] (8.5, 6.5);
	\draw (8, 2.5) to[capacitor, /tikz/circuitikz/bipoles/length=0.700cm, l_={$C_{GN}$}] (8.5, 2.5);
	\draw (4.25, 4.25) -- (5, 4.25);
	\draw (4.25, 4.75) -- (5, 4.75);
	\draw (3.75, 4.75) -- (3.5, 4.75);
	\draw (3.75, 4.25) -- (3.5, 4.25);
\end{tikzpicture}
\end{center}

da questo vediamo le diverse capacità da considerare:
\begin{itemize}
	\item La capacità del filo che collega la prima porta alla seconda, che chiamiamo $C_W$. Questa è la capacità di una traccia nel caso dei circuiti integrati, o di un filo conduttore vero e proprio, ecc..
	\item La capacità fra drain e body dei MOSFET, che chiamiamo $C_{DB}$;
	\item La capacità fra gate e drain dei MOSFET, che chiamiamo $C_{GD}$;
	\item Le capacità gra gate e body dei MOSFET, che chiamiamo $C_{GP}$ per il PMOS (con body a $V_{DD}$) e $C_{GN}$ per l'NMOS (con body a ground).
\end{itemize}

Vediamo di calcolare la capacità equivalente al nodo drain della prima porta. Visto che la maggior parte di queste (tutte tranne le $G_{CD}$)sono in parallelo:
$$
C_{EQ} = C_W + C_{DBN} + C_{DBP} + C_{GN} + C_{GP} + 2 C_{GDN} + 2 C_{GDP}
$$

Possiamo considerare la $C_{DBP}$ e la $C_{GP}$ come in parallelo, in quanto alle variazioni i nodi a $V_{DD}$ sono fissi (e noi stiamo calcolando la risposta ai transitori).

Una particolarità è data dalle $G_{CD}$, in quanto queste effettivamente sono poste fra $V_{in}$ e $V_{out}$:
$$
V_{GD} = V_{in} - V_{out}
$$
Possiamo allora calcolare la variazione di carica sulla capacità in un transitorio, come:
$$
\Delta Q = C_{GD} \cdot 2 \Delta V
$$
dove $\Delta V$ è la variazione di tensione su $V_{in}$. Assumendo linearità, La $\Delta V$ sulla $V_{in}$ è uguale alla $-\Delta V$ sulla $V_{out}$, per cui la variazione di tensione totale è:
$$
\Delta V_{in} - \Delta V_{out} = \Delta V_{in} - (-\Delta V_{out}) = 2 \Delta V
$$
da cui il $2$.

Possiamo quindi prendere l'intera capacità come una singola collegata alla $V_{out}$ e a ground, e di capacità $2 C_{GD}$. Chiamiamo questo fenomeno anche effetto \textit{Miller}.

\par\smallskip

Calcoliamo quindi la potenza dissipata nei due transitori.
\begin{itemize}
\item Transizione \textit{alto $\rightarrow$ basso}: abbiamo che in questo caso la tensione iniziale di uscita è $V_{out} = V_{DD}$, mentre la tensione finale è $0 \, \mathrm{V}$. La potenza immagazzinata all'inizio della transizione dalla capacità equivalente sarà:
$$
E_H = \frac{1}{2} C_{EQ} V_{DD}^2 
$$
mentre l'energia finale sarà:
$$
E_L = 0
$$
per cui l'energia totale dissipata sarà:
$$
E_{HL} = E_H - E_L = \frac{1}{2} C_{EQ} V_{DD}^2
$$
Questa energia sarà quella dissipata sul NMOS in conduzione in fase di scarica, che ricordiamo a $V_{in}$ alta entra in conduzione (e si scalda per effetto Joule quando la capacità ci scarica sopra).

Notiamo che propriamente si rovrebbe calcolare:
$$
\int_0^\tau V_C(t) i_C(t) \, dt
$$
sulla capcità, cioè applicare $P = VI$. Visto che non conosciamo con esattezza questi valori, però, possiamo applicare il bilancio energetico a fine transizione;

\item Transizione \textit{basso $\rightarrow$ alto}: abbiamo che in questo caso la potenza dissipata è quella usata dal PMOS per caricare la $C_{EQ}$. Anche qui applichiamo il principio di conservazione dell'energia, e calcoliamo il bilancio energetico a fine transizione.

	Iniziamo dalla potenza erogata dalla $V_{DD}$ (cioè dall'alimentatore che mantiene la $V_{DD}$ al suo valore) durante la transizione. Ricordando che la tensione iniziale di uscita è $V_{out} = 0 \, \mathrm{V}$, e la tensione finale è $V_{DD}$. L'energia erogata è quindi:
$$
E_{DD} = \int_0^\tau V_{DD} i_{DD}(t) \, dt = V_DD Q
$$
dove $Q$ è la carica del condensatore portato a $V_{DD}$:
$$
Q = V_{DD} C_{EQ}
$$
per cui l'energia erogata diventa:
$$
E_{DD} = V_{DD}^2 C_{EQ}
$$

La potenza immagazzinata nella capacità equivalente a fine transizione sarà sempre la solita:
$$
E_{H} = \frac{1}{2} C_{EQ} V_{DD^2}
$$

Abbiamo quindi che l'energia totale dissipata sarà data da quella erogata dalla batteria meno quella immagazzinata dalla capacità (che non viene persa):
$$
E_{LH} = E_{DD} - E_H = \frac{1}{2} C_{EQ} V_{DD}^2
$$

\end{itemize}

Abbiamo trovato un risultato che ha senso fisicamente: l'energia che usiamo per portare il CMOS in uno stato stabile è uguale, che si transisca da alto a basso o da basso a alto.

\par\smallskip

Per passare da considerazioni energetiche a considerazioni di potenza assumiamo di commutare la porta CMOS un certo numero di volte su unità di tempo, cioè ad una certa \textit{frequenza}. Visto che abbiamo 2 commutazioni per ciclo, la potenza dissipata risulterà:
$$
P_D = f C_{EQ} V_{DD}^2
$$

Questa equazione è fondamentale e rappresenta il fatto che la potenza che dissipiamo va:
\begin{itemize}
	\item Linearmente con la frequenza, per cui clock a frequenza maggiore portano ad una maggiore dissipazione di potenza;
	\item Linearmente con la capacità equivalente $C_{EQ}$. Storicamente, questa è stata ridotta diminuendo la dimensione fisica dei componenti (anche se vediamo che questo non si può fare all'infinito);
	\item Quadraticamente con la tensione, per cui dimezzando l'alimentazione si riduce di quattro volte la potenza dissipata. Questo giustifica ad esempio come mai si preferisce la logica a $3.3 \, \mathrm{V}$ anziché a $5 \, \mathrm{V}$;
\end{itemize}

\subsection{Sintesi di porte logiche CMOS}
Come abbiamo anticipato, la sintesi di porte logiche più complicate si fa estendendo la struttura dell'inverter. Si vanno quindi a formare 2 reti:
\begin{itemize}
	\item Una rete di \textit{pull-up}, collegata fra l'alimentazione $V_{DD}$ e l'uscita $V_{out}$;
	\item Una rete di \textit{pull-down}, collegata fra l'uscita $V_{out}$ e il ground.
\end{itemize}

Le variabili di ingresso vengono collegate contemporaneamente alle reti di pull-up e di pull-down. Quando la rete di pull up è in conduzione, quella di pull down è disattiva (e quindi la $V_{out}$ è alta), mentre viceversa quando la rete di pull down è in conduzione, quella di pull up è disattiva (e quindi la $V_{out}$ è bassa).

\subsubsection{Configurazioni di MOSFET}
Vediamo alcune primitive che possiamo costruire usando i transistori MOS.
Di base vale la seguente regola:
\begin{itemize}
	\item I MOS in serie formano degli "AND";
	\item I MOS in parallelo formano degli "OR".
\end{itemize}

Questo risulta chiaro guardando i circuiti, ad esempio nel caso degli NMOS:
\begin{center}
\begin{tikzpicture}[transform shape]
	% Paths, nodes and wires:
	\node[nmos] at (0, 0.77){};
	\node[nmos] at (0, -0.77){};
	\draw (-0.98, 0.77) -| (-1.5, 1);
	\draw (-0.98, -0.77) -| (-1.5, -0.5);
	\node[shape=rectangle, minimum width=0.465cm, minimum height=0.465cm] at (-1.5, 1.25){} node[anchor=south, align=center, text width=0.077cm, inner sep=6pt] at (-1.5, 1){$V_A$};
	\node[shape=rectangle, minimum width=0.465cm, minimum height=0.465cm] at (-1.5, -0.25){} node[anchor=south, align=center, text width=0.077cm, inner sep=6pt] at (-1.5, -0.5){$V_B$};
	\draw[-stealth] (0, 1.54) -| (0, 2) -- (1, 2);
	\node[shape=rectangle, minimum width=0.465cm, minimum height=0.465cm] at (1.25, 2){} node[anchor=west, align=left, text width=0.077cm, inner sep=6pt] at (1, 2){$V_{out}$};
	\node[tlground] at (0, -1.54){};
	\node[nmos] at (4, -0.27){};
	\node[nmos] at (6, -0.31){};
	\draw (4, 0.5) -| (4, 0.96) -- (6, 0.96) -| (6, 0.46);
	\node[shape=rectangle, minimum width=0.465cm, minimum height=0.465cm] at (2.75, -0.29){} node[anchor=east, align=right, text width=0.077cm, inner sep=6pt] at (2.8, -0.29){$V_A$};
	\node[shape=rectangle, minimum width=0.465cm, minimum height=0.465cm] at (4.77, -0.31){} node[anchor=east, align=right, text width=0.077cm, inner sep=6pt] at (4.8, -0.31){$V_B$};
	\node[tlground] at (4, -1.04){};
	\node[tlground] at (6, -1.08){};
	\draw[-stealth] (6, 0.96) -- (7, 0.96);
	\node[shape=rectangle, minimum width=0.465cm, minimum height=0.465cm] at (7.25, 0.96){} node[anchor=west, align=left, text width=0.077cm, inner sep=6pt] at (7, 0.96){$V_{out}$};
\end{tikzpicture}
\end{center}

\begin{itemize}
	\item A sinistra abbiamo la parte NMOS della porta NAND, in quanto entrambi i MOSFET devono condurre (e quindi sia $V_A$ che $V_B$ devono essere alti) perché $V_{out}$ sia collegato a massa. Abbiamo quindi implementato l'operazione:
$$
\overline{Y} = A \cdot B
$$
chiamando $Y$ l'uscita.
	\item A destra abbiamo la parte NMOS della porta NOR, in quanto basta che uno dei MOSFET conduca (cioè $V_A$ o $V_B$ sia alto) perché $V_{out}$ sia collegato a massa. Abbiamo quindi implementato l'operazione:
$$
\overline{Y} = A + B
$$
\end{itemize}

La stessa cosa, ma ad ingressi invertiti, si ha nel caso dei PMOS:
\begin{center}
\begin{tikzpicture}[transform shape]
	% Paths, nodes and wires:
	\node[nmos] at (0, 0.77){};
	\node[nmos] at (0, -0.77){};
	\draw (-0.98, 0.77) -| (-1.5, 1);
	\draw (-0.98, -0.77) -| (-1.5, -0.5);
	\node[shape=rectangle, minimum width=0.465cm, minimum height=0.465cm] at (-1.5, 1.25){} node[anchor=south, align=center, text width=0.077cm, inner sep=6pt] at (-1.5, 1){$V_A$};
	\node[shape=rectangle, minimum width=0.465cm, minimum height=0.465cm] at (-1.5, -0.25){} node[anchor=south, align=center, text width=0.077cm, inner sep=6pt] at (-1.5, -0.5){$V_B$};
	\draw[-stealth] (0, -1.54) -- (0, -2.04) -- (1, -2.04);
	\node[shape=rectangle, minimum width=0.465cm, minimum height=0.465cm] at (1.25, -2.04){} node[anchor=west, align=left, text width=0.077cm, inner sep=6pt] at (1, -2.04){$V_{out}$};
	\node[nmos] at (4, -0.27){};
	\node[nmos] at (6, -0.27){};
	\node[shape=rectangle, minimum width=0.465cm, minimum height=0.465cm] at (2.75, -0.29){} node[anchor=east, align=right, text width=0.077cm, inner sep=6pt] at (2.8, -0.29){$V_A$};
	\node[shape=rectangle, minimum width=0.465cm, minimum height=0.465cm] at (4.77, -0.27){} node[anchor=east, align=right, text width=0.077cm, inner sep=6pt] at (4.8, -0.27){$V_B$};
	\draw[-stealth] (6, -2) -- (7, -2);
	\node[shape=rectangle, minimum width=0.465cm, minimum height=0.465cm] at (7.25, -2){} node[anchor=west, align=left, text width=0.077cm, inner sep=6pt] at (7, -2){$V_{out}$};
	\draw (4, -1.04) -| (4, -2) -- (6, -2) |- (6, -1);
	\node[vcc] at (4, 0.5){};
	\node[vcc] at (6, 0.5){};
\end{tikzpicture}
\end{center}

\begin{itemize}
	\item A sinistra abbiamo la parte PMOS della porta AND ad ingressi negati, in quanto entrambi i MOSFET devono condurre (e quindi sia $V_A$ che $V_B$ devono essere bassi) perché $V_{out}$ sia collegato alla $V_{DD}$. Abbiamo quindi implementato l'operazione:
$$
Y = \overline{A} \cdot \overline{B}
$$
\item A destra abbiamo la parte PMOS della porta OR ad ingressi negati, in quanto basta che uno dei MOSFET conduca (cioè $V_A$ o $V_B$ sia bassa) perché $V_{out}$ sia collegato alla $V_{DD}$. Abbiamo quindi implementato l'operazione:
$$
Y = \overline{A} + \overline{B}
$$
\end{itemize}

\subsubsection{Implementazione della NOR}
Iniziamo ad implementare alcune semplici porte.

\par\smallskip

\noindent
\textbf{\textsf{Porta NOR}}

\noindent
Vediamo un caso semplice, cioè quello della NOR a 2 ingressi. Per ricavarla dobbiamo innanzitutto prendere l'espressione logica che vogliamo implementare:
$$
Y = \overline{A + B}
$$
e quindi ricavare le espressioni della PUN e della PDN prendendo rispettivamente entrate e uscita negate: 
\begin{itemize}
	\item PUN: $Y = \overline{A} \cdot \overline{B}$
	\item PDN: $\overline{Y} = A + B$
\end{itemize}

Abbiamo già introdotto queste reti, per cui possiamo sintetizzare il circuito come segue:
\begin{center}
\begin{tikzpicture}[transform shape]
	% Paths, nodes and wires:
	\node[pmos] at (0, 0.77){};
	\node[pmos] at (0, 2.31){};
	\node[nmos] at (-1, -0.77){};
	\node[nmos] at (1, -0.77){};
	\draw (-1, -0) -- (0, -0) -- (1, -0);
	\node[tlground] at (-1, -2){};
	\node[tlground] at (1, -2){};
	\node[vcc] at (0, 3.08){};
	\draw (-2.26, -0.77) -| (-3, 2.25);
	\draw (-0.98, 2.31) -| (-3, 2.25);
	\node[jump crossing] at (-1, -1.68){};
	\draw (-1, -1.82) |- (-1, -2);
	\draw (1, -2) |- (1, -1.54);
	\draw (0.02, -0.77) -| (-0.5, -1);
	\draw (-0.5, -1) |- (-0.86, -1.68);
	\node[jump crossing, rotate=90] at (-2.12, -0.77){};
	\draw (-1.14, -1.68) -| (-2.12, -0.91);
	\draw (-2.12, -0.63) |- (-0.98, 0.77);
	\draw[-stealth] (1, -0) -- (2, -0);
	\draw (-2.12, -1.68) -| (-4, -1.5);
	\draw (-3, 2.31) -| (-4, 2.5);
	\node[shape=rectangle, minimum width=0.465cm, minimum height=0.465cm] at (2.25, -0){} node[anchor=west, align=left, text width=0.077cm, inner sep=6pt] at (2, -0){$V_{out}$};
	\node[shape=rectangle, minimum width=0.465cm, minimum height=0.465cm] at (-4, 2.75){} node[anchor=south, align=center, text width=0.077cm, inner sep=6pt] at (-4, 2.5){$V_A$};
	\node[shape=rectangle, minimum width=0.465cm, minimum height=0.465cm] at (-4, -1.25){} node[anchor=south, align=center, text width=0.077cm, inner sep=6pt] at (-4, -1.5){$V_B$};
\end{tikzpicture}
\end{center}

Notiamo quindi che una particolarità delle porte in tecnologia CMOS è che il funzionamento della porta è assicurato quando le reti PUN e PDN sono duali. Questa non è una condizione necessaria, però, e vedremo che porte (come la XOR) non la rispettano. 

\end{document}
